\documentclass[12pt,a4paper]{article}
\usepackage{fontspec}\usepackage{polyglossia}\setdefaultlanguage{ukrainian}
\setmainfont{NewComputerModern10}\setsansfont{Noto Sans}\setmonofont{Noto Sans Mono}[Scale=.82]
\usepackage{amsmath,amssymb,booktabs,caption,enumitem,float,geometry,graphicx,hyperref,microtype,titlesec,xcolor}
\geometry{left=2cm,right=2cm,top=2cm,bottom=2cm}\setlength{\parindent}{1.25cm}\linespread{1.08}
\titleformat{\section}{\normalfont\fontsize{14}{16}\selectfont\bfseries}{\thesection}{.6em}{}\titlespacing*{\section}{0pt}{1.3em}{.6em}
\captionsetup{font=small,labelfont=bf,labelsep=endash}
\usepackage[newfloat]{minted}\SetupFloatingEnvironment{listing}{name=Лістинг}
\setminted{fontsize=\footnotesize,linenos,breaklines,autogobble,frame=single,framesep=2mm,bgcolor=gray!5,tabsize=4}
\hypersetup{unicode,colorlinks=true,linkcolor=black,pdftitle={ІО-41, Давидчук А. М., Лабораторна робота №7},pdfauthor={Давидчук А. М.}}
\newcommand{\resultDir}{../results}\newcommand{\studentGroup}{ІО-41}\newcommand{\studentName}{Давидчук А. М.}\newcommand{\reportYear}{2026}
\begin{document}\pagenumbering{gobble}
\begin{titlepage}
\begin{center}\includegraphics[width=.7\textwidth]{../../../../../../templates/letterhead.png}\end{center}
\begin{center}\large Національний технічний університет України\\ «Київський політехнічний інститут імені Ігоря Сікорського»\\[1em]Факультет інформатики та обчислювальної техніки\\Кафедра обчислювальної техніки\end{center}
\vfill\begin{center}\textbf{\LARGE Методи та технології штучного інтелекту}\\[2em]\textbf{\Large Лабораторна робота №7}\\[.7em]\parbox{.9\textwidth}{\centering «Знаходження мінімуму та максимуму функцій за допомогою генетичних алгоритмів»}\end{center}
\vfill\begin{flushright}\begin{minipage}{.48\textwidth}Виконав: студент 3 курсу ФІОТ,\\групи \studentGroup\\\textit{\studentName}\\[1em]Перевірив: \textit{Польшакова О. М.}\end{minipage}\end{flushright}
\vfill\begin{center}Київ --- \reportYear\end{center}
\end{titlepage}\pagenumbering{arabic}

\section{Тема та мета роботи}
\noindent\textbf{Тема:} пошук екстремумів функцій генетичним алгоритмом.\\
\noindent\textbf{Мета:} реалізувати генетичний алгоритм із дійсним кодуванням, знайти мінімум функції однієї змінної й максимум функції двох змінних, дослідити збіжність та різноманітність популяції.

\section{Варіант і постановка задачі}
У методичному файлі згадано додаток Б, але його немає. Тому використано функції варіанта 6 з наявної таблиці додатка А, яка вже застосовувалась у попередніх роботах:
\begin{align}
f(x)&=\sin(|x|)+\cos(|x|), & x&\in[-4;1],\\
z(x,y)&=x\sin(y), & x,y&\in[-1;3].
\end{align}
Потрібно мінімізувати $f(x)$ і максимізувати $z(x,y)$. Відповідно до методички розмір популяції дорівнює 10, а кількість поколінь --- 100.

\section{Теоретичні відомості}
Генетичний алгоритм --- стохастичний метод глобального пошуку, у якому кожне можливе рішення подається хромосомою. Популяція містить кілька хромосом, а функція пристосованості визначає якість кожної з них. Для задачі максимізації пристосованість збігається зі значенням функції, а для мінімізації використовується протилежний знак.

Основний цикл складається з оцінювання, відбору, кросоверу, мутації та формування нового покоління. Відбір надає перевагу якісним рішенням; кросовер поєднує координати двох батьків; мутація додає нові варіанти і не дозволяє популяції передчасно стати однаковою. Елітизм без змін переносить найкращу хромосому в наступне покоління, тому знайдений результат не втрачається.

У роботі застосовано дійсне кодування: хромосома для першої задачі містить $[x]$, а для другої --- $[x,y]$. Турнірний відбір обирає найкращого з трьох випадкових кандидатів. Арифметичний кросовер формує нащадка
\begin{equation}c=\alpha p_1+(1-\alpha)p_2.\end{equation}
Мутація додає гаусове відхилення з імовірністю 0,20 для кожної координати. Її масштаб зменшується протягом еволюції: спочатку алгоритм досліджує простір, наприкінці --- уточнює рішення.

\section{Програмна реалізація}
Обидві цільові функції передаються до одного універсального оптимізатора.
\begin{listing}[H]
\inputminted[firstline=18,lastline=25]{python}{../src/main.py}
\caption{Функції варіанта 6}
\end{listing}

Основний цикл генетичного алгоритму реалізує елітизм, два турніри для вибору батьків, кросовер, мутацію та обмеження координат допустимим інтервалом.
\begin{listing}[H]
\inputminted[firstline=28,lastline=64]{python}{../src/main.py}
\caption{Універсальний генетичний алгоритм}
\end{listing}

\section{Результати}
\begin{table}[H]\centering
\begin{tabular}{lrrr}\toprule Задача & $x$ & $y$ & Значення функції\\\midrule
Мінімум $f(x)$ & $-3{,}92698890$ & --- & $-1{,}41421356$\\
Максимум $z(x,y)$ & $3{,}00000000$ & $1{,}57080537$ & $3{,}00000000$\\\bottomrule
\end{tabular}\caption{Знайдені екстремуми}\end{table}

Для першої функції теоретичний мінімум $-\sqrt2$ досягається, коли $|x|=5\pi/4\approx3{,}92699$. Знайдене значення збігається з ним до наведених знаків. Для другої задачі максимум дорівнює 3, бо найбільше допустиме $x=3$, а $\sin(y)=1$ при $y=\pi/2\approx1{,}5708$.

\begin{figure}[H]\centering\includegraphics[width=.96\textwidth]{\resultDir/convergence.png}\caption{Зміна найкращого та середнього значень у поколіннях}\end{figure}
У першій задачі мінімум знайдено приблизно до п'ятого покоління, у другій --- до четвертого. Подальші покоління уточнюють координати і стабілізують середню якість популяції. Швидка збіжність пояснюється невеликою розмірністю функцій.

\begin{figure}[H]\centering\includegraphics[width=.72\textwidth]{\resultDir/minimum_1d.png}\caption{Графік функції однієї змінної та знайдений мінімум}\end{figure}
Червона точка розташована в найнижчій області графіка в заданому інтервалі, що візуально підтверджує правильність результату.

\begin{figure}[H]\centering\includegraphics[width=.8\textwidth]{\resultDir/maximum_2d.png}\caption{Поверхня функції двох змінних та знайдений максимум}\end{figure}
Максимум лежить на межі $x=3$ і біля вершини синусоїди за координатою $y$. Це узгоджується з аналітичними властивостями функції.

\begin{figure}[H]\centering\includegraphics[width=.75\textwidth]{\resultDir/diversity.png}\caption{Зміна різноманітності популяцій}\end{figure}
На початку координати особин розкидані по всьому інтервалу. Після відбору їхнє стандартне відхилення зменшується, але мутація періодично створює нові варіанти й підтримує глобальний пошук.

Оскільки генетичний алгоритм є стохастичним, додатково виконано 10 незалежних запусків з різними початковими значеннями генератора випадкових чисел. Середнє значення знайденого мінімуму становило $-1{,}4142135252$ зі стандартним відхиленням $4{,}07\cdot10^{-8}$, а середнє значення максимуму --- $2{,}9999997112$ зі стандартним відхиленням $7{,}66\cdot10^{-7}$.

\begin{figure}[H]\centering\includegraphics[width=.86\textwidth]{\resultDir/repeated_runs.png}\caption{Похибка екстремумів у 10 незалежних запусках}\end{figure}
На графіку показано абсолютне відхилення від аналітичних значень $-\sqrt2$ та 3 у логарифмічному масштабі. В усіх запусках похибка є малою, отже результат не залежить від одного випадково вдалого початкового набору особин.

\section{Висновки}
Реалізовано генетичний алгоритм із дійсним кодуванням, турнірним відбором, арифметичним кросовером, адаптивною мутацією та елітизмом. За популяції 10 особин і 100 поколінь знайдено мінімум $-1{,}41421356$ у точці $x=-3{,}92698890$ та максимум 3 у точці $(3;1{,}57080537)$. Десять незалежних запусків дали стандартні відхилення результатів менш як $10^{-6}$, тому реалізація не лише збігається з аналітичними екстремумами, а й демонструє стабільну роботу.

\section{Контрольні питання}
\begin{enumerate}[leftmargin=*]
\item \textit{Що таке генетичний алгоритм і на якій ідеї він базується?}\\
Це еволюційний метод пошуку, який імітує природний відбір, успадкування, схрещення та мутації.
\item \textit{Які основні етапи генетичного алгоритму?}\\
Ініціалізація популяції, оцінювання, відбір, кросовер, мутація, формування нового покоління та перевірка зупинки.
\item \textit{Що таке популяція?}\\
Це множина кандидатів-розв'язків, які одночасно досліджуються на поточному поколінні.
\item \textit{Яка роль функції пристосованості?}\\
Вона перетворює якість рішення на числову оцінку, за якою виконуються відбір і порівняння особин.
\item \textit{Як відбувається відбір?}\\
Більш пристосовані особини мають більшу ймовірність стати батьками; у роботі використано турнір із трьох кандидатів.
\item \textit{Що таке кросовер?}\\
Це створення нащадка шляхом поєднання генетичної інформації двох батьків.
\item \textit{Які типи кросоверу застосовують?}\\
Одноточковий, багатоточковий, рівномірний, арифметичний та blend-кросовер для дійсних координат.
\item \textit{У чому полягає мутація?}\\
Вона випадково змінює окремі гени, додаючи нові рішення і допомагаючи вийти з локального екстремуму.
\item \textit{Як підтримують різноманітність популяції?}\\
Мутацією, обмеженням сили відбору, випадковим додаванням особин та спеціальними штрафами за схожість.
\item \textit{Які критерії зупинки використовують?}\\
Максимальну кількість поколінь або часу, досягнення потрібної якості чи відсутність покращення протягом заданого числа поколінь.
\item \textit{Для яких задач застосовують генетичні алгоритми?}\\
Для складної глобальної, дискретної, комбінаторної та багатокритеріальної оптимізації.
\item \textit{Які переваги ГА над класичними методами?}\\
Вони не потребують похідних, працюють із розривними функціями та одночасно досліджують кілька областей простору.
\item \textit{Які недоліки мають генетичні алгоритми?}\\
Вони виконують багато оцінок функції, не гарантують точного оптимуму і чутливі до налаштування параметрів.
\item \textit{Як параметри впливають на результат?}\\
Велика популяція покращує покриття простору, мутація підтримує різноманітність, а кількість поколінь визначає доступний час пошуку.
\item \textit{Де генетичні алгоритми застосовують на практиці?}\\
У логістиці, проєктуванні, плануванні, машинному навчанні, біоінформатиці, економіці та оптимізації мереж.
\end{enumerate}
\end{document}
